\documentclass[11pt, a4paper]{moderncv}
\usepackage[margin=.7in]{geometry}
\usepackage[french]{babel}

%-------------------------------------------------------------------------------
%                Infos
%-------------------------------------------------------------------------------
\name{Dorian}{Boully}
\born{17 mars 2000}
\address{50 Rue de la Goutte d'Or, 75018, Paris, France}{}{}
\phone[mobile]{(+33)6~41~68~08~93}
\email{dorian.boully@polytechnique.edu}
\title{}

%-------------------------------------------------------------------------------
%                Style customization
%-------------------------------------------------------------------------------
\moderncvstyle{classic}
\moderncvcolor{burgundy}
\moderncvhead{1}%
\usepackage[utf8]{inputenc}
\usepackage[T1]{fontenc}
\usepackage{lmodern}

\renewcommand*{\mobilephonesymbol}{}
\renewcommand*{\bornsymbol}{}
\renewcommand*{\emailsymbol}{}

% Lists customizations
\usepackage[]{enumitem}
\setlist{itemsep=0pt, parsep=0pt}
\setlist[itemize]{label=\raisebox{0.4ex}{\tiny$\bullet$}, topsep=0pt,partopsep=0.5\baselineskip}

%-------------------------------------------------------------------------------
%                Document
%-------------------------------------------------------------------------------
\begin{document}
\makecvtitle%
\begin{center}
    \vspace{-2em}
  \quotestyle{École Polytechnique (4\ieme{} année) -- M2 Mathématiques
  fondamentales IMJ-PRG}
  \vspace{1.5em}
\end{center}

\section{Formation académique}
\cventry{depuis 2024}{M2 Mathématiques fondamentales}{IMJ-PRG}{}{}{\normalsize
  \begin{itemize}
    \item En cours : Convergence de spectres, Classes caractéristiques ;
    \item Cours validés : Surfaces de Riemann, Topologie algébrique des variétés, Algèbre homologique, Schémas ;
    \item Exposé au séminaire $\infty$-catégories sur les (co)limites ;
    \item Exposé au groupe de travail Géométrie hyperkählérienne sur la conjecture de Calabi.
  \end{itemize}
}

\cventry{depuis 2020}{École Polytechnique}{Parcours d'approfondissement en mathématiques}{}
{}{\normalsize
    Projets avec mémoire :
    \begin{itemize}
        \item Calcul de Malliavin. Supervision : Anne-Sophie de Suzzoni ;
        \item Représentations de groupes et équation de Dirac. Supervision : A-S. de Suzzoni et S. Munier ;
        \item Trous noirs colorés. Supervision : Cécile Huneau et Lucas Chesnel ;
        \item Polynômes lorentziens et conjecture de Mason. Supervision : Omid Amini.
    \end{itemize}}

\cventry{2017 -- 2020}{CPGE MPSI/MP}{Lycée Louis-le-Grand}{}{}{}

\vspace{0.5em}
\section{Expériences professionnelles et stages}
\cventry{février 2026 --}{Stage de recherche M2}{IMJ-PRG}{}{}
{Topologie des variétés algébriques réelles.\\
Tuteur : Ilia Itenberg.}
\cventry{avril -- août 2023}{Stage de recherche (M1/3A de l'X)}{IMJ-PRG}{}{}
{Introduction à la géométrie kählérienne (théorème de Kodaira).\\
Tutrice : Eleonora Di Nezza.}
\cventry{2022 -- 2024}{Colleur de mathématiques (MP2I)}{Lycée Saint-Louis}{}{}{}
\cventry{juin -- août 2022}{Stagiaire analyste quantitatif}{Pointsbet Europe}{Dublin}{}
{\normalsize Modèle de cotation de paris sportifs in-game pour la NFL.}
\cventry{décembre 2020 -- avril 2021}{Stage militaire}{25\ieme{} RGA}{BA 702, Avord, France}
{}
{\normalsize Étude de faisabilité pour la qualification des terrains sommaires A400M.}

\vspace{0.5em}
\section{Projets personnels}
\cvitem{janvier 2026}{Rédaction d'une preuve détaillée du théorème d'uniformisation de Riemann.}
\cvitem{janvier 2026}{Rédaction de notes pour le cours Convergences de spectres
et notes fondamentales de Nalini Anantharaman au Collège de France.}
\vspace{0.5em}
\section{Langues}
\begin{itemize}
\item \textbf{Français :} langue maternelle
\item \textbf{Anglais :}  C1 certifié TOEFL
\item \textbf{Espagnol :} B2 certifié DELE
\end{itemize}

\section{Centres d'intérêt}
\cvitem{Sport}{Badminton (niveau national).}

\end{document}